Here we will introduce $\perpp$-ideals, which will be used in relation to partial actions of inverse semigroups in the next chapter. As usual, we fix a locally compact Hausdorff space $X$, a continuous function $\theta$ from $X$ to a Hausdorff space $H$ and a weakly regular family $\mathcal{A}\subseteq C_c(X,\theta)$.

\begin{definition}\label{definitionstrongcover}\index{Cover!strong}
A finite family $B\subseteq\mathcal{A}$ is said to be a \emph{strong cover} of an element $a\in\mathcal{A}$ if there is another finite family $\widetilde{B}\subseteq\mathcal{A}$ such that:
\begin{enumerate}
\item[(i)] For all $\widetilde{b}\in\widetilde{B}$, there is some $b\in B$ with $\widetilde{b}\Subset b$;
\item[(ii)] $\widetilde{B}$ is a cover of $a$ (Definition \ref{definitiondisjointcover}).
\end{enumerate}
\end{definition}

The following proposition shows that strong covers encode information about the closures of sets. It is a direct consequence of the definition of $\Subset$ and Lemma \ref{lemmacovers}.

\begin{proposition}\label{propositiongoodformofstrongcover}
A finite family $B\subseteq\mathcal{A}$ is a strong cover of $a\in\mathcal{A}$ if and only if $\supp(a)\subseteq\bigcup_{b\in B}\sigma(b)$.
\end{proposition}

\begin{definition}\index{Ideal!$\perpp$-ideal}\label{definitionperppideal}
A \emph{$\perpp$-ideal} in $\mathcal{A}$ is a subset $I\subseteq A$ such that, for all $a\in\mathcal{A}$,
\[a\in I\iff \text{there is a finite subset }B\subseteq I\text{ which is a strong cover of }a.\]
\end{definition}

\begin{example}
Suppose $\mathcal{A}$ is a vector sublattice of $C_c(X,\mathbb{R})$ or $C(X,\mathbb{C})$ (self-adjoint in the latter case), where $X$ is a locally compact Hausdorff space. Then a subset $I\subseteq\mathcal{A}$ is a $\perpp$-ideal if and only if $I$ is an additive subgroup and
\[f\in I\iff\exists g\in I\text{ such that }f\Subset g.\ntag\label{propertyofidealinC(X)}\]

Indeed, first suppose $I$ is a $\perpp$-ideal. If $f\in I$, then there is a finite strong cover $B\subseteq I$ of $f$. For every $b\in B$, take another finite strong cover $C(b)\subseteq I$ of $b$. Then $\bigcup_{b\in B}C(b)$ is a strong cover of $g=\sum_{b\in B}|\operatorname{Re}(b)|+|\operatorname{Im}(b)|$, thus $g\in I$ and $f\Subset g$, which proves implication ``$\Rightarrow$'' of (\ref{propertyofidealinC(X)}). Implication ``$\Leftarrow$'' is immediate since $f\Subset g$ if and only if $\left\{g\right\}$ is a strong cover of $f$.

To prove that $I$ is additive, let $f_1,f_2\in I$. The implication ``$\Rightarrow$'' of (\ref{propertyofidealinC(X)}), yields $g_1,g_2\in I$ with $f_i\Subset g_i$. Then $\left\{g_1,g_2\right\}$ is a finite strong cover of $f_1+f_2$, so $f_1+f_2\in I$.

Now suppose that $I$ is additive and satisfies property (\ref{propertyofidealinC(X)}), and let us show that $I$ is a $\perpp$-ideal. Of course, if $f\in I$ then $f$ admits a finite cover by a single element of $I$ (by the ``$\Rightarrow$'' part of (\ref{propertyofidealinC(X)})). Conversely, suppose $f$ admits a finite strong cover $B\subseteq I$. From (\ref{propertyofidealinC(X)}) take, for all $b\in B$, $g_b\in I$ with $b\Subset g_b$, so that $\widetilde{b}:=|\operatorname{Re}(b)|+|\operatorname{Im}(b)|\Subset g_b$ as well. This in turn implies $\widetilde{b}\in I$, so $g=\sum_{b\in B}\widetilde{b}\in I$ and $f\Subset g$, so we conclude that $f\in I$ again from (\ref{propertyofidealinC(X)}).
\end{example}

We will now prove that the lattice of open subsets of a space $X$ is order-isomorphic to the lattice of $\perpp$-ideals of a weakly regular tuple $(X,\theta,\mathcal{A})$.

\begin{definition}\nomenclature[Z]{$\mathbf{I}(U)$}{$\perpp$-ideal associated to an open set $U$}\nomenclature[Z]{$\mathbf{U}(I)$}{Open set associated to a $\perpp$-ideal $I$}
Suppose $(X,\theta,\mathcal{A})$ is weakly regular. Given an open set $U\subseteq X$, denote $\mathbf{I}(U)=\left\{f\in\mathcal{A}:\supp(f)\subseteq U\right\}$. Given a $\perpp$-ideal $I\subseteq\mathcal{A}$, denote $\mathbf{U}(I)=\bigcup_{f\in I}\sigma(f)$.
\end{definition}

Note that $\mathbf{I}(U)$ is a $\perpp$-ideal of $\mathcal{A}$: If $f\in\mathbf{I}(U)$ then we use weak regularity of $\mathcal{A}$ and compactness of $\supp(f)$ to find a finite collection $B\subseteq\mathcal{A}$ such that
\[\supp(f)\subseteq\bigcup_{b\in B}\sigma(b)\subseteq U\]
In particular, $B\subseteq\mathbf{I}(U)$ and $B$ is a strong cover of $f$. This provides the implication ``$\Rightarrow$'' of Definition \ref{definitionperppideal}, whereas ``$\Leftarrow$'' is trivial.

\begin{theorem}\label{theoremperppideals}
Suppose $(X,\theta,\mathcal{A})$ is weakly regular.
\begin{enumerate}
\item[(a)] For every $\perpp$-ideal $I$ of $\mathcal{A}$, $I=\mathbf{I}(\mathbf{U}(I))$;
\item[(b)] For every open subset $U\subseteq X$, $U=\mathbf{U}(\mathbf{I}(U))$;
\item[(c)] The map $U\mapsto\mathbf{I}(U)$ is an order isomorphism between the lattices of open sets of $X$ and $\perpp$-ideals of $\mathcal{A}$.
\end{enumerate}
\end{theorem}
\begin{proof}
\begin{enumerate}
\item[(a)] Let $I$ be a $\perpp$-ideal. The inclusion $I\subseteq \mathbf{I}(\mathbf{U}(I))$ follows easily from the property of $\perpp$-ideals: if $f\in I$, then take a finite strong cover $B\subseteq I$ of $f$, so that $\supp(f)\subseteq\bigcup_{b\in B}\sigma(b)\subseteq\mathbf{U}(I)$.

Conversely, if $f\in\mathbf{I}(\mathbf{U}(I))$ we use compactness of $\supp(f)$ to find a finite family $B\subseteq I$ with $\supp(f)\subseteq\bigcup_{b\in B}\sigma(b)$, so $B$ is a strong cover of $f$ and therefore $f\in I$.
\item[(b)] Suppose $U\subseteq X$ is open. By weak regularity of $(X,\theta,\mathcal{A})$, we have
\[U=\bigcup_{\substack{f\in\mathcal{A}\\\supp(f)\subseteq U}}\sigma(f)=\bigcup_{f\in\mathbf{I}(U)}\sigma(f)=\mathbf{U}(\mathbf{I}(U)).\]
\item[(c)] The previous items prove that $U\mapsto\mathbf{I}(U)$ is a bijection, with inverse $I\mapsto\mathbf{U}(I)$. It is clear that both maps are order-preserving.\qedhere
\end{enumerate}
\end{proof}

Assume $(X,\theta,\mathcal{A})$ is weakly regular. Let $\widehat{\mathcal{A}}$ be the collection of maximal $\perpp$-ideals of $\mathcal{A}$, and endow it with the topology generated by the sets
\[U(f)=\left\{I\in\widehat{\mathcal{A}}:\exists g\Subset f\text{ such that }g\not\in I\right\},\qquad f\in\mathcal{A}.\]
By Theorem \ref{theoremperppideals}, we obtain a bijection $\xi:X\to\widehat{\mathcal{A}}$, $\xi(x)=\mathbf{I}(X\setminus\left\{x\right\})$. Since for all $x\in X$ and $f\in\mathcal{A}$,
\[x\in\sigma(f)\iff\exists g\Subset f\text{ such that } x\in\supp(g)\iff\xi(x)\in U(f),\]
then $\xi(\sigma(f))=U(f)$, which proves that $\xi$ is a homeomorphism. This provides another proof of Theorem \ref{maintheorem}.